\clearpage
\fancyhead[L]{\small WIRE SPACINGS.}
\hypertarget{native-450-fadenintervalle}{}
\begin{center}
{\large WIRE SPACINGS.}
\end{center}

\begin{center}
\rule{0.5\linewidth}{0.4pt}\\[2pt]
Astronomische Nachrichten. Vol. II. No. 43. Pages 371\,..\,376. November 1823.\\[2pt]
\rule{0.5\linewidth}{0.4pt}
\end{center}

\begin{center}
A new method of determining the distances between the wires in meridian telescopes.
\end{center}

The only method that seems to be in use for determining wire spacings in meridian telescopes is to observe the transits of Polaris or of other stars near the pole
\SourceLeaf{451}%
whose declinations are known. It combines convenience with great precision when it can be used under favourable conditions. In seasons when unfavourable weather tends to prevail, one often has to wait a long time and make many unsuccessful attempts before one can observe even a single complete transit. When, as so often happens, the stars quiver markedly as they pass through the field of view instead of moving steadily, the method can yield only much lower precision. These considerations would lose much of their importance if the final determination of the spacings could always be postponed until the conditions for applying that method were exceptionally favourable. But one would like to be able to carry out early the reductions that require the wire spacings, and delay can also be risky, especially since a reticle of delicate spider threads is so easily damaged. I have sometimes had to replace the reticle repeatedly at short intervals because a thread had gone slack or broken by accident, even before a single complete transit of Polaris could be observed in good seeing. One cannot count on the spacings being identical, even if the new threads are laid in the same previously scribed grooves, as experience has often taught me. These reasons prompted me to look for a new way of determining the wire spacings that could be used to advantage alongside the usual method and in its place under unfavourable conditions. The one I recently tried seems to serve the purpose.

Just as all mutually parallel rays striking the objective meet at a single point in the plane where the reticle is to be placed, so conversely all the rays that travel in the opposite direction from a single point in this plane and strike the objective become parallel after passing through it. Rays that start from different points, by contrast, emerge from the objective at angles to one another exactly equal to the separation of those points as expressed in angular terms when the instrument is used. Thus if the eyepiece end of the telescope faces the sky or another bright surface, a farsighted eye looking through the objective would see the reticle clearly and at the proper separation, in the sense just mentioned, if the eye were sensitive enough to see objects so fine. What the unaided eye cannot do becomes possible with a second telescope whose eyepiece is set as the eye requires for very distant objects. In fact the threads set at the focus of a \textsc{Frauenhofer} objective can be seen in this way with the greatest clarity and definition.
The wire spacings can now also be measured this way with great precision. I used the same 12-inch theodolite by \textsc{Ertel} with which I observed the angles in the measurement of a meridian arc. Its own reticle consists of one horizontal wire and two parallel wires half a minute apart. In a new eyepiece that \textsc{Repsold} made for this purpose these wires are astonishingly fine. The measurement comes out very neatly when the wires of the meridian instrument, seen through both objectives, are centred between the parallel wires.

It goes without saying that for this operation the eyepiece of the meridian telescope must be removed entirely, or at least replaced by one of lower magnification through which the wires whose spacings are to be measured can be seen at the same time.

The distance of the theodolite from the objective does not matter, but one will prefer not to make it too large, to avoid receiving some rays too close to the edge of the objective or even with only half illumination.

\Needspace{3\baselineskip}
\SourceLeaf{452}%
A properly levelled theodolite gives the angles projected onto the horizon rather than the actual angles. The values measured directly therefore still have to be reduced to the true angles. For distances as small as those in my meridian instruments, multiplication by the cosine of the inclination is sufficient for this purpose. A meridian circle itself gives that inclination with the greatest precision. With a simple transit telescope that gives the inclination to a minute at best, one can just as well obtain it from the small altitude circle of the theodolite. In that case it is advisable to avoid too great an inclination.

The following measurements should be regarded only as an initial trial. They were by no means made with all the precautions it will be advisable to take. The theodolite and its wooden tripod stood on the false floor of planks between the piers of the meridian instruments and were not independent of the observer’s chair. Its position was somewhat awkward at the considerable inclination, and the light was not especially good for taking the readings. If these difficulties are avoided, there is all the more reason to expect closer agreement. Each figure is the result of 10 repetitions.

\Needspace{23\baselineskip}
\emph{Distances from the middle wire in the telescope of the \textup{\textsc{Reichenbach}} meridian circle to the 6 other wires.}

\begin{center}
1823 Oct.\ 23. \quad Inclination $=23^\circ\,58'\,52''$

\begingroup\renewcommand{\arraystretch}{1.1}
\begin{tabular}{r@{\;}c@{}|r|r}
 & & Horizontal & True dist. \\
\cline{3-4}
\smash{\raisebox{-0.6em}{I}} & \smash{\raisebox{-0.6em}{$\bigg\{$}} & $11'\,25''050$ & $10'\,25''93$ \\
 & & $11\quad23.000$ & $10\quad24.86$ \\
\smash{\raisebox{-0.6em}{II}} & \smash{\raisebox{-0.6em}{$\bigg\{$}} & $7\quad45.900$ & $7\quad5.69$ \\
 & & $7\quad46.275$ & $7\quad6.03$ \\
\smash{\raisebox{-0.6em}{III}} & \smash{\raisebox{-0.6em}{$\bigg\{$}} & $3\quad58.225$ & $3\quad37.66$ \\
 & & $3\quad57.750$ & $3\quad37.23$ \\
\smash{\raisebox{-0.6em}{V}} & \smash{\raisebox{-0.6em}{$\bigg\{$}} & $3\quad58.850$ & $3\quad38.24$ \\
 & & $3\quad59.375$ & $3\quad38.71$ \\
\smash{\raisebox{-0.6em}{VI}} & \smash{\raisebox{-0.6em}{$\bigg\{$}} & $7\quad53.625$ & $7\quad12.74$ \\
 & & $7\quad53.875$ & $7\quad12.98$ \\
\smash{\raisebox{-0.6em}{VII}} & \smash{\raisebox{-0.6em}{$\bigg\{$}} & $11\quad17.600$ & $10\quad19.11$ \\
 & & $11\quad17.775$ & $10\quad19.28$ \\
\end{tabular}
\endgroup

\medskip
Nov.\ 6. \quad Inclination $=23^\circ\,57'\,30''$

\begingroup\renewcommand{\arraystretch}{1.1}
\begin{tabular}{r@{\;}c@{}|r|r}
\smash{\raisebox{-0.6em}{I}} & \smash{\raisebox{-0.6em}{$\bigg\{$}} & $11'\,24''800$ & $10'\,25''29$ \\
 & & $11\quad25.575$ & $10\quad26.51$ \\
\end{tabular}
\endgroup
\end{center}

The wire spacings were also calculated from three observed transits of Polaris, only one of which could be observed in full. The seeing was neither especially good nor especially poor. The first two were lower transits in daylight and the third was an upper transit at night.

\begin{center}
\begingroup\small\setlength{\tabcolsep}{4pt}
\begin{tabular}{l|r|r|r|r|r|r}
 & I & II & III & V & VI & VII \\
Oct.\ 23 & $10'\,26''91$ & .\,.\,.\,. & $3'\,38''54$ & $3'\,38''12$ & $7'\,11''27$ & $10'\,19''31$ \\
Nov.\ 5 & $10\quad26.90$ & $7'\,7''01$ & $3\quad38.33$ & $3\quad39.39$ & $7\quad10.41$ & $10\quad19.50$ \\
\phantom{Nov.\ }6 & $10\quad24.31$ & .\,.\,.\,. & $3\quad35.12$ & $3\quad40.22$ & .\,.\,.\,. & .\,.\,.\,. \\
\end{tabular}
\endgroup
\end{center}

If one applies the principles developed in my Theoria Combinationis Observationum to the determinations made with the theodolite, the mean error of a wire spacing from 10 repetitions is $= 0''67$. Applying them to the Polaris transits, on the other hand, gives the mean error from one transit as $= 1''38$. The former determination would accordingly have twice the precision or four times the weight. However, the number of measurements
\SourceLeaf{453}%
on each side is too small for the mean error to be inferred from them with great reliability. For Polaris it applies in any case only to similar seeing, since it would turn out considerably smaller in exceptionally good seeing. In general, however, the usefulness of the new procedure is at least sufficiently confirmed by this.

Another no less interesting application of these principles is to determine the magnification of telescopes. I will only touch on it here. If the eyepiece is set for a farsighted eye, objects viewed through the telescope in reverse appear reduced by the same factor as they are magnified in ordinary use. They are seen so clearly that they can readily be magnified again through a second telescope. One can therefore measure by repetition with the theodolite the angle between two objects viewed in this way and compare it with the angle between the directions from the centre of the eyepiece to the objects themselves. The objects need not be very far away. With fairly strong magnification, a few feet may suffice. For greater precision, however, when the angles at the eyepiece are fairly large, due account must be taken of the fact that the quantities proportional to the magnifications are, more precisely, the tangents of the angles the directions make with the optical axis, rather than the angles themselves. It is therefore easiest to choose the objects so that the direction of the optical axis lies very nearly midway between the directions from the centre of the eyepiece to the two objects. One then compares the tangent of half the angle at the eyepiece with the tangent of half the angle obtained from the theodolite observation after proper reduction to the true value. A heliometer can also be used in place of the theodolite and gives the true value directly.
This procedure is capable of far greater precision than the others that have been used for the same purpose.
